\section{Orbit stability: further proofs and results}\label{sec:sup-stability}
We supplement the stability results of Appendix~\ref{sec:fqe-orbit-stability} of the paper and Section~\ref{sec:fqe-stability-proofs}
with a constructive distance certificate, the proofs of the two sharp critical-locus theorems,
a one-parameter refinement on $P=1$, and the proof of the sharp entropy coefficient on the critical locus.
Throughout, $\mathcal M$ is the union of the two HS local-unitary orbits,
including global phase, and
\[
 d(\psi)=\min_{\chi\in\mathcal M}\|\psi-\chi\|,\qquad
 \Phi(\psi)=\max_{\chi\in\mathcal M}|\langle\psi,\chi\rangle|^2,\qquad
 d^2=2(1-\sqrt{\Phi}).
\]
For unit states, $Q=F+(P-1)$ is the sextic moment defect and
$\delta=E(M_4)-E$ is the entropy deficit, as in the paper.

\subsection{A direct certificate from three Bell-pair products}\label{sec:sup-bellframe}
Let $\tilde w=\epsilon^{\otimes4}\bar w$ be the spin flip, as in the paper; in the magic matrix coordinates it is ordinary
conjugation (Section~\ref{sec:sup-general-envelope}). Hence a unit state $w$ with
$\langle\tilde w,w\rangle=0$ has the form $(x+iy)/\sqrt2$ for orthonormal real vectors $x,y$ in the magic coordinates; put
$V=\operatorname{span}_{\mathbb R}(x,y)$. A \emph{Bell-pair product} is a product of two maximally entangled Bell states with arbitrary local
unitaries on the four qubits, with its pairing specified; it is real up to phase in these coordinates.

\begin{proposition}[Three-Bell-frame certificate]\label{prop:fqe-bell-frame}
Let $w$ be unit with $\langle\tilde w,w\rangle=0$. Choose one unit
Bell-pair product $b_j$ in each of the three pairings, and put
\[
 \varepsilon_j=1-2|\langle b_j,w\rangle|^2,\qquad
 E_B=\varepsilon_1+\varepsilon_2+\varepsilon_3.
\]
Then $\varepsilon_j\ge0$ and
\[
 \boxed{\quad d(w)^2\le\frac92 E_B.\quad}
\]
The vectors $b_j$ need not maximise their overlaps. The proof constructs
at most six candidate HS states, two for each unordered pair of pairings.
\end{proposition}
\begin{proof}
Choose phases so the $b_j$ are real. Their squared distances to $V$ are
exactly $\varepsilon_j$. Bell-pair products in distinct pairings have
$|\langle b_i,b_j\rangle|\le1/2$: in the magic matrix coordinates for
the first pairing they are a real rank-one matrix and $O/2$, $O\in O(4)$.
The proper and improper components describe the two other pairings.
This description follows by expanding the standard Bell products and
using $SU(2)\otimes SU(2)\to SO(4)$.

Let $G$ be the $3\times3$ real Gram matrix of the $b_j$, and let
$\lambda_3$ be its smallest eigenvalue. The rank-two projection on $V$
gives $\lambda_3\le E_B$. Since $\operatorname{tr}G=3$,
\[
 \operatorname{tr}G^2\ge\frac{(3-\lambda_3)^2}{2}+\lambda_3^2
                     \ge\frac92-3E_B.
\]
Consequently, for $\Delta_{ij}=1/4-|\langle b_i,b_j\rangle|^2$,
\begin{equation}\label{eq:fqe-bell-gram}
 \Delta_{ij}\ge0,\qquad \sum_{i<j}\Delta_{ij}\le\frac32 E_B.
\end{equation}

Fix a pair $(i,j)$ and use real rotations to make $b_i=E_{00}$
and $b_j=O/2$. Adjust its real sign so
$c=O_{00}/2=|\langle b_i,b_j\rangle|\in[0,1/2]$.
The lower $3\times3$ block of $O$ has singular values $1,1,2c$:
its product with its transpose is $I-zz^T$, with $\|z\|^2=1-4c^2$.
Choose its polar factor $R\in O(3)$, including a factor on any nullspace,
and put $O'=\operatorname{diag}(1,R)$, $b'_j=O'/2$.
The real plane $S=\operatorname{span}(b_i,b'_j)$ is spanned by
$E_{00}$ and $\operatorname{diag}(0,R)/\sqrt3$. Its two normalised
complex-null rays are precisely HS states, since their magic matrices
are $\operatorname{diag}(1/\sqrt2,\,\pm iR/\sqrt6)$.

Put $S_0=\operatorname{span}(b_i,b_j)$. Its Gram matrix has
smallest eigenvalue $1-c$.
Consequently its leakage from $V$ satisfies
\[
 \ell_{ij}:=\operatorname{tr}(P_{S_0}(I-P_V))
 \le\frac{\varepsilon_i+\varepsilon_j}{1-c}
 \le2(\varepsilon_i+\varepsilon_j).
\]
The planes $S_0,S$ share $b_i$. Their other principal-angle cosine is
\[
 r_c=\frac{1+c}{\sqrt{3(1-c^2)}}
     =\sqrt{\frac{1+c}{3(1-c)}}.
\]
Indeed the polar-aligned lower block of $O$ has trace $2+2c$.
The distance squared between the nearer normalised null rays of these
planes is $1-r_c$. Moreover
\[
 g(c):=1-r_c\le\frac74(1/4-c^2)=\frac74\Delta_{ij}.
\]
To check the last inequality, square its nonnegative right-hand
comparison with $r_c$. After multiplication by $3(1-c)$ the difference
is
$(1-2c)[13+c(525-462c)+588c^3(1-2c)]/256\ge0$.

For any two real two-planes with principal-angle cosines $s_1,s_2$,
choose orthonormal bases in which their overlap matrix has singular values
$s_1,s_2$. Reversing an orientation replaces a null ray
$(x+iy)/\sqrt2$ by its conjugate. Thus the largest overlap between
their normalised null rays is $(s_1+s_2)/2$, and their minimum
squared distance (after a phase choice) is
$2-s_1-s_2\le2-s_1^2-s_2^2$, the projection leakage.
The triangle inequality therefore gives
$d(w)\le\sqrt{\ell_{ij}}+\sqrt{g(c)}$.
Average its square over the three unordered pairs, use
\eqref{eq:fqe-bell-gram}, and apply Cauchy--Schwarz to the cross term:
\[
 d(w)^2\le\left(\sqrt{\frac43}+\sqrt{\frac78}\right)^2E_B
          \le\frac92 E_B,
\]
since $(\sqrt{4/3}+\sqrt{7/8})^2<9/2$. The last inequality is strict when $E_B>0$.
\end{proof}

For a general unit state choose phase and real orthonormal $x,y$ so
$\psi=\sqrt{(1+t)/2}\,x+i\sqrt{(1-t)/2}\,y$,
where $t=|\langle\tilde\psi,\psi\rangle|$, and set $w=(x+iy)/\sqrt2$.
At $t=1$, choose any real unit $y$ orthogonal to $x$.
The same constructive certificate gives
\[
 d(\psi)\le\sqrt{2-\sqrt{1+t}-\sqrt{1-t}}+\sqrt{9E_B(w)/2}.
\]
This individual-state certificate requires only chosen Bell products and
linear algebra; the entropy-only bound is proved in
Section~\ref{sec:fqe-global-overlap}.

% The proofs of the paper's two sharp critical-locus theorems (Section sec:fqe-critical-orbit-stability).
\subsection{The critical locus: signed Birkhoff reduction and real-density model}
\label{sec:fqe-critical-orbit-proof}
A state is critical when all four one-qubit marginals are $I/2$. On this locus, both the overlap comparison and its Hilbert-space distance
version have optimal constants.

\begin{restated}[Sharp critical-state orbit stability; from Theorem~\ref{thm:fqe-global-overlap}, CO-1--CO-3]{Theorem}{thm:fqe-critical-orbit-sharp}
For every critical unit pure four-qubit state, with $\Phi,d,Q,k$
as in Appendix~\ref{sec:fqe-orbit-stability},
\begin{align*}
 Q(\psi)&\ge\max\{1-\Phi(\psi),\;3/2-2\Phi(\psi)\},\\
 Q(\psi)&\ge\frac{2+\sqrt2}{4}d(\psi)^2,\\
 \delta(\psi)&\ge k\bigl(\max\{1-\Phi,3/2-2\Phi\}+5\tau/2\bigr)
       \ge k\,\frac{2+\sqrt2}{4}\,d(\psi)^2.
\end{align*}
The coefficient one in $Q\ge1-\Phi$, and the coefficient
$(2+\sqrt2)/4$ in \eqref{eq:fqe-critical-distance}, are both optimal:
Bell-pair products attain them. The comparison $Q\ge1-\Phi$ holds for all states
(Theorem~\ref{thm:fqe-global-overlap}).
If the critical state also has $\tau=0$, equivalently if $P=1$, then
\begin{equation*}
 \Phi\ge\frac{2+\sqrt3}{6},\qquad
 Q\ge\frac3{4-\sqrt3}(1-\Phi)\ge\frac{3+\sqrt3}{4}d^2.
\end{equation*}
These coefficients and the fidelity floor are optimal, attained by cluster states.
\end{restated}
Section~\ref{sec:sup-nullsharp} below proves a one-parameter sharpening of
\eqref{eq:fqe-null-critical-sharp} in the Bloch coordinates of \eqref{eq:fqe-null-critical-bloch}. The optimal entropy
coefficient on the whole critical locus is larger: $\delta\ge C_*d^2$ with $C_*=0.19057894701\ldots$ \eqref{eq:fqe-p1-cstar}
(Theorem~\ref{thm:fqe-p1-sharp}), optimal by its equality states; $k(2+\sqrt2)/4$ in \eqref{eq:fqe-critical-entropy}
is the weaker coefficient of the moment route.

We first remove the continuous orbit optimisation exactly.
For complex diagonal matrices $D_v,D_w$ of any common size,
\begin{equation}\label{eq:fqe-signed-diagonal-overlap}
 \max_{A,B\in O(m)}|\operatorname{tr}(D_v^\dagger A D_wB^T)|
 =\max_{\pi,\,\varepsilon_i=\pm1}
                  \left|\sum_i\bar v_i\varepsilon_iw_{\pi(i)}\right|.
\end{equation}
Indeed $|A_{ij}B_{ij}|$ is doubly substochastic, by Cauchy--Schwarz
on each row and column. Complete it to a doubly stochastic matrix;
for any common phase the real overlap is bounded by the largest
assignment sum of the absolute real coefficient entries. Birkhoff's
permutation decomposition gives this bound, and signed permutation
matrices attain it. Maximising the phase proves
\eqref{eq:fqe-signed-diagonal-overlap}.

The magic normal form makes a critical state $D_v$. Here and below the coordinates of $v\in\mathbb C^4$, and the
rows and columns of $4\times4$ matrices, are indexed $0,\dots,3$: $v_i$ is the coefficient of $m_{i+1}\otimes m_{i+1}$ in
$G(v)$ of Appendix~\ref{sec:fqe-normal-form} of the paper, so $D_v=\operatorname{diag}(v)$.
Up to global phase, the union $\mathcal M$ is the $O(4)\times O(4)$ orbit
of $Z_*$: a common orientation reversal in both factors fixes $Z_*$
up to proper rotations, whereas a single reversal gives its conjugate
up to proper rotations. Thus its overlap with $D_v$ is exactly the
maximum over the following 32 diagonal representatives:
\[
 \mathcal H=\left\{\eta:\ \eta_i=1/\sqrt2,\quad
         \eta_j=\varepsilon_j i/\sqrt6\ (j\ne i),\quad i=0,1,2,3\right\}.
\]
Its exact spin-flip-dependent covering profile is proved in
Lemma~\ref{lem:fqe-critical-real-density} below.

The real cut transforms are $H_0=I$ and
\[
 H_1=\frac12\begin{pmatrix}
 1&-1&1&-1\\-1&1&1&-1\\1&1&1&1\\-1&-1&1&1
 \end{pmatrix},\qquad
 H_2=\frac12\begin{pmatrix}
 1&-1&-1&-1\\-1&1&-1&-1\\1&1&-1&1\\-1&-1&-1&1
 \end{pmatrix}.
\]
With $\Delta=\operatorname{diag}(1,-1,1,-1)$, so that $\Delta v$ are the coordinates $(a,b,c,d)$ of $G_{abcd}$, these are
the transforms $H_\pm$ of Section~\ref{sec:sup-sextic-zeros} conjugated by $\Delta$: $H_1=\Delta H_+\Delta$, and $H_2$
is $\Delta H_-\Delta$ with its rows permuted and re-signed, which only permutes the squared moduli $|e_i^TH_2v|^2$ of the cut spectrum.

\begin{proof}[Proof of Theorem~\ref{thm:fqe-critical-orbit-sharp}]
The comparison $Q(\psi)\ge1-\Phi(\psi)$ holds for all states by
Theorem~\ref{thm:fqe-global-overlap}. When $\Phi(\psi)\ge\tfrac12$,
$1-\Phi(\psi)\ge\tfrac32-2\Phi(\psi)$, so \eqref{eq:fqe-critical-fidelity}
holds. Assume henceforth that $\Phi(v)\le\tfrac12$ for a unit critical
normal vector $v\in\mathbb C^4$.

The three real orthogonal cut transforms $H_0=I,H_1,H_2$ displayed above
preserve $N=\|v\|^2=1$, $\tau=|\sum_i v_i^2|^2$ and $Q(v)$, and permute
the $32$ representatives $\mathcal H$ up to the sixth-root-of-unity
phases $(\pm1\pm i\sqrt3)/2$: for any $\eta\in\mathcal H$, one coordinate
is $1/\sqrt2$ and three are $\pm i/\sqrt6$, so under $H_1$ or $H_2$ the
row whose relative signs on those three coordinates match $\eta$ has inner
product $\frac{\pm1\pm i\sqrt3}{2\sqrt2}$ and the other three rows have
$\frac{\pm1\pm i\sqrt3}{2}(\mp i/\sqrt6)$ (CO-1). Thus $H_0,H_1,H_2$,
coordinate permutations, global phase and coordinate sign flips preserve
both $Q(v)$ and $\Phi(v)$. Replace $v$ by $H_mv$ ($m\in\{0,1,2\}$) and
permute coordinates so that
$p_0:=|v_0|^2=\max_{m,i}|e_i^TH_mv|^2$ is the largest of all twelve cut
probabilities. Choose global phase so that $v_0=r_0>0$, and flip signs
$v_j\mapsto\pm v_j$ ($1\le j\le3$) so that $y_j:=\operatorname{Im}v_j\ge0$,
writing $v_j=x_j+iy_j$ and $p_j=x_j^2+y_j^2\le p_0=r_0^2$. Put
$X_2:=\sum_{j=1}^3x_j^2$, $Y_2:=\sum_{j=1}^3y_j^2=1-p_0-X_2$,
$S_x:=\sum_{j=1}^3|x_j|$ and $S_y:=\sum_{j=1}^3y_j$.

Testing $\Phi(v)$ against
$Z_*=(1/\sqrt2,i/\sqrt6,i/\sqrt6,i/\sqrt6)^T\in\mathcal H$ gives
\[
 \Phi(v)\ge\Phi_{\rm im}(v):=\left(\frac{r_0}{\sqrt2}+\frac{S_y}{\sqrt6}\right)^2
 =\frac12-\frac12X_2+\frac13\Bigl(\sqrt3\,r_0S_y-Y_2+\sum_{1\le k<\ell\le3}y_ky_\ell\Bigr).
\]
Since $\Phi_{\rm im}(v)\le\Phi(v)\le\tfrac12$, we have
$S_y\le\sqrt3(1-r_0)$, hence $Y_2\le S_y^2\le3(1-r_0)^2$ and
$X_2=1-r_0^2-Y_2\ge2(1-r_0)(2r_0-1)$. Moreover, since
$0\le y_j\le\sqrt{p_j}\le r_0$ and $\sqrt3<\tfrac{26}{15}<\tfrac74$
($26^2=676>675$), we also have
$Y_2\le r_0S_y\le\tfrac{26}{15}r_0(1-r_0)\le\tfrac74r_0(1-r_0)$ and
$X_2\ge\tfrac1{15}(1-r_0)(15-11r_0)$.

Up to sign, the eight rows of $H_1$ and $H_2$ are
$\tfrac12(1,\varepsilon_1,\varepsilon_2,\varepsilon_3)$ for
$\varepsilon\in\{\pm1\}^3$, so the eight cut probabilities in $H_1v$ and
$H_2v$ are
\[
 \lambda_\varepsilon
 =\frac14\left[\left(r_0+\sum_{j=1}^3\varepsilon_jx_j\right)^2
   +\left(\sum_{j=1}^3\varepsilon_jy_j\right)^2\right]
 =\frac14\left(1+2\sum_{j=1}^3\varepsilon_jc_{0j}
   +2\sum_{1\le k<\ell\le3}\varepsilon_k\varepsilon_\ell c_{k\ell}\right)\le p_0,
\]
where $c_{ij}:=\operatorname{Re}(\bar v_iv_j)=x_ix_j+y_iy_j$
($c_{0j}=r_0x_j$). Choosing $\varepsilon_j=\operatorname{sgn}(x_j)$ gives
$(r_0+S_x)^2\le4p_0$, so $S_x\le r_0$. For each $j\in\{1,2,3\}$, averaging
$\lambda_\varepsilon\le p_0$ over the four sign vectors with
$\varepsilon_j=\operatorname{sgn}(x_j)$ cancels the other terms and yields
$\tfrac14(1+2r_0|x_j|)\le r_0^2$, so $|x_j|\le\frac{4r_0^2-1}{2r_0}$ and
$X_2\le(\max_j|x_j|)S_x\le2r_0^2-\tfrac12$. Combining this with
$Y_2\le\tfrac74r_0(1-r_0)$ gives
\[
 1=r_0^2+X_2+Y_2\le3r_0^2-\frac12+\frac74r_0(1-r_0),\qquad
 0\le\frac14(5r_0-3)(r_0+2),
\]
hence $r_0\ge\tfrac35$. Furthermore, for $\varepsilon_j=\operatorname{sgn}(x_j)$
we have $(\sum_{j=1}^3\varepsilon_jy_j)^2\ge Y_2-2\sum_{k<\ell}y_ky_\ell
=2Y_2-S_y^2\ge2(1-p_0-X_2)-3(1-r_0)^2$, so
\[
 (r_0+S_x)^2\le4p_0-2(1-p_0-X_2)+3(1-r_0)^2=(3r_0-1)^2+2X_2
 \le\left(3r_0-1+\frac{X_2}{3r_0-1}\right)^2,
\]
and therefore $S_x\le\min\bigl(r_0,\,2r_0-1+\frac{X_2}{3r_0-1}\bigr)$. For
$s_2:=\sum_{1\le k<\ell\le3}|x_k||x_\ell|=\tfrac12(S_x^2-X_2)$, if
$X_2\le(1-r_0)(3r_0-1)$ then
\[
 2s_2\le\left(2r_0-1+\frac{X_2}{3r_0-1}\right)^2-X_2
 =(2r_0-1)^2-\frac{X_2\bigl((1-r_0)(3r_0-1)-X_2\bigr)}{(3r_0-1)^2}
 \le(2r_0-1)^2,
\]
while if $X_2\ge(1-r_0)(3r_0-1)$ then
$2s_2\le r_0^2-X_2\le(2r_0-1)^2$. Thus $2s_2\le(2r_0-1)^2$ and
$q_4:=\sum_{1\le k<\ell\le3}x_k^2x_\ell^2\le s_2^2\le\tfrac14(2r_0-1)^4$
for all $X_2\ge0$.

On the critical locus, expanding $Q(v)$ in terms of $p_i=|v_i|^2$ and
$c_{ij}=\operatorname{Re}(\bar v_iv_j)$ gives the exact identity (CO-1)
\begin{equation}\label{eq:fqe-critical-three-term}
 Q(v)=\frac12\Bigl(\sum_{i=0}^3p_i\Bigr)^3
      +18\sum_{j=1}^3D_j\bigl(c_{0j}^2-c_{k\ell}^2\bigr),
\end{equation}
where for each $j\in\{1,2,3\}$ and $\{k,\ell\}=\{1,2,3\}\setminus\{j\}$ we
set $D_j:=p_0+p_j-p_k-p_\ell=2(p_0+p_j)-1=1-2(p_k+p_\ell)\le
D_{\max}:=\min(1,4p_0-1)$. Using $c_{0j}^2=p_0x_j^2$,
$c_{k\ell}^2\le2(x_k^2x_\ell^2+y_k^2y_\ell^2)$ when $D_j\ge0$, and
$c_{k\ell}^2\ge0$ when $D_j<0$, we obtain with $D_j^+:=\max(D_j,0)$ the
decoupled lower bound
\[
 Q(v)+2\Phi(v)-\frac32\ge Q(v)+2\Phi_{\rm im}(v)-\frac32
 \ge\mathcal Y(r_0,y)+\mathcal X(p;x^2),
\]
where
\begin{align*}
 \mathcal Y(r_0,y)&:=\frac23\Bigl(\sqrt3\,r_0S_y-Y_2
   +\sum_{1\le k<\ell\le3}y_ky_\ell\Bigr)-36\sum_{j=1}^3D_j^+y_k^2y_\ell^2,\\
 \mathcal X(p;x^2)&:=\sum_{j=1}^3(18p_0D_j-1)x_j^2
   -36\sum_{j=1}^3D_j^+x_k^2x_\ell^2.
\end{align*}

We show that both $\mathcal Y(r_0,y)$ and $\mathcal X(p;x^2)$ are
nonnegative (CO-2). For $\mathcal Y(r_0,y)$, the four-term AM--GM
inequality on $(2r_0+1)+3(2r_0-1)+4(1-r_0)+4(1-r_0)=6$ gives
$48(4r_0^2-1)(1-r_0)^2\le(3/2)^4=81/16$, so
$9D_{\max}S_y^2\le27(4r_0^2-1)(1-r_0)^2\le729/256<3$. Since
$Y_2=S_y^2-2\sum_{k<\ell}y_ky_\ell$ and
$\sum_{k<\ell}y_k^2y_\ell^2\le\tfrac14S_y^2\sum_{k<\ell}y_ky_\ell$
(as $y_ky_\ell\le\tfrac14(y_k+y_\ell)^2\le\tfrac14S_y^2$), we have
\[
 \mathcal Y(r_0,y)\ge\frac23S_y(\sqrt3\,r_0-S_y)
   +\bigl(2-9D_{\max}S_y^2\bigr)\sum_{1\le k<\ell\le3}y_ky_\ell
 \ge\frac23S_y(\sqrt3\,r_0-S_y)-\sum_{1\le k<\ell\le3}y_ky_\ell.
\]
Using $\sum_{k<\ell}y_ky_\ell\le\tfrac13S_y^2$ and
$S_y\le\sqrt3(1-r_0)$ gives
$\mathcal Y(r_0,y)\ge\tfrac13S_y(2\sqrt3\,r_0-3S_y)
\ge\tfrac{\sqrt3}{3}S_y(5r_0-3)\ge0$.

For $\mathcal X(p;x^2)$, substituting
$D_j=(2p_0-1)+2x_j^2+2y_j^2\ge(2p_0-1)+2x_j^2$ and
$\sum_{j=1}^3x_j^4=X_2^2-2q_4$ into the first sum yields
\[
 \mathcal X(p;x^2)\ge\bigl(18p_0(2p_0-1)-1\bigr)X_2+36p_0(X_2^2-2q_4)
   -36\sum_{j=1}^3D_j^+x_k^2x_\ell^2.
\]
If $X_2=0$, then $\mathcal X=0$; assume $X_2>0$. On
$r_0\in[\tfrac{19}{25},1]$, we have $p_0\ge\tfrac{361}{625}>\tfrac12$, so
$D_j=1-2(p_k+p_\ell)\ge0$ and $D_j\le1-2(x_k^2+x_\ell^2)$. Hence
$\sum_{j=1}^3D_j^+x_k^2x_\ell^2\le q_4-2(X_2q_4-3x_1^2x_2^2x_3^2)
\le(1-\tfrac43X_2)q_4$ by $3x_1^2x_2^2x_3^2\le\tfrac13X_2q_4$, and since
$2p_0+1-\tfrac43X_2>0$ and $q_4\le\tfrac13X_2^2$,
\[
\begin{aligned}
 \frac{\mathcal X(p;x^2)}{X_2}&\ge18p_0(2p_0-1)-1-12(1-p_0)X_2+16X_2^2\\
 &=\left(4X_2-\frac32(1-p_0)\right)^2+\frac{135p_0^2-54p_0-13}{4}
 \ge\frac{16573}{78125}>0.
\end{aligned}
\]
On $r_0\in[\tfrac35,\tfrac{19}{25}]$, using $D_j^+\le D_{\max}$ and
$q_4\le\tfrac14(2r_0-1)^4$ gives
\[
 \frac{\mathcal X(p;x^2)}{X_2}\ge18p_0(2p_0-1)-1+36p_0X_2
   -\frac{9(2p_0+D_{\max})(2r_0-1)^4}{X_2},
\]
which is strictly increasing in $X_2>0$. On $[\tfrac35,\tfrac{16}{25}]$,
setting $D_{\max}=4p_0-1$ and
$X_2\ge X_{2,1}(r_0):=\tfrac1{15}(1-r_0)(15-11r_0)$ and multiplying by
$15X_{2,1}(r_0)>0$ gives the degree-six polynomial
$P_1(r_0)=15\bigl((36r_0^4-18r_0^2-1)X_{2,1}(r_0)+36r_0^2X_{2,1}(r_0)^2
-9(6r_0^2-1)(2r_0-1)^4\bigr)$, which is strictly concave on
$[\tfrac35,1]$ ($P_1''(\tfrac35+u)=-368208u^4-\tfrac{2057376}{5}u^3
-\tfrac{2990232}{25}u^2-\tfrac{1591704}{125}u-\tfrac{226078}{625}<0$ for all
$u\ge0$) with $P_1(\tfrac35)=\tfrac{3693}{78125}>0$ and
$P_1(\tfrac{16}{25})=\tfrac{216609912}{1220703125}>0$. On
$[\tfrac{16}{25},\tfrac{19}{25}]$, setting $D_{\max}=1$ and
$X_2\ge X_{2,2}(r_0):=2(1-r_0)(2r_0-1)$ and multiplying by
$\frac{X_{2,2}(r_0)}{2r_0-1}=2(1-r_0)>0$ gives the degree-five polynomial
$P_2(r_0)=72r_0^5-432r_0^4+432r_0^3-54r_0^2-52r_0+7$, which is strictly
concave on $[\tfrac35,1]$
($P_2''(\tfrac35+u)=-36\bigl(8u^2(9-5u)+\tfrac{288}{5}u+3\bigr)<0$ for
$u\in[0,\tfrac25]$) with $P_2(\tfrac{16}{25})=\tfrac{988047}{9765625}>0$
and $P_2(\tfrac{19}{25})=\tfrac{570453}{9765625}>0$. Thus
$\mathcal X(p;x^2)\ge0$ in all cases, proving \eqref{eq:fqe-critical-fidelity}.

Let $x=\sqrt{\Phi}$ and $c=(2+\sqrt2)/4$. The $32$ overlaps
$q_\eta=|\langle\eta,v\rangle|^2$ sum to $8$, since the frame operator is $8I$ (CO-1), so
$\Phi=\max_\eta q_\eta\ge1/4$ and $x\ge1/2$.
For $x\ge1/\sqrt2$, the first bound gives
$1-x^2\ge2c(1-x)$. For $1/2\le x\le1/\sqrt2$, the second gives
\[
 3/2-2x^2-2c(1-x)
 =2(1/\sqrt2-x)\bigl(x-(2-\sqrt2)/4\bigr)\ge0.
\]
Since $d^2=2(1-x)$, this proves \eqref{eq:fqe-critical-distance}.
Equation~\eqref{eq:fqe-orbit-surrogate} gives \eqref{eq:fqe-critical-entropy}, with
$\mathcal D=Q+5\tau/2$ on critical states and the bound just proved as its last step. Bell products have
$Q=\Phi=1/2$, hence $d^2=2-\sqrt2$; both constants are sharp.
The null-critical refinement is proved in the next paragraph.
\end{proof}

\paragraph{Two auxiliary Bloch spheres on the null-critical locus.}
Let $M=(m_1,m_2,m_3,m_4)$ be the unitary $4\times4$ matrix whose columns are the magic basis vectors of
Appendix~\ref{sec:fqe-normal-form} of the paper, written in the basis
$|00\rangle,|01\rangle,|10\rangle,|11\rangle$, so that $Mv=\sum_kv_km_{k+1}$ for $v\in\mathbb C^4$. For the critical
normal form $D_v$, regard $Mv$ as an auxiliary two-qubit vector.
Its coefficient matrix has determinant $\sum_jv_j^2/2$.
Thus $\tau=0$ is equivalent to $Mv=a\otimes b$ for unit auxiliary
qubits $a,b$. These are coordinates on the critical normal form, not
physical tensor factors of the original four-qubit state.
Let $r,s\in S^2$ be their Bloch vectors. Direct Bell probabilities give
\begin{equation}\label{eq:fqe-null-critical-bloch}
 Q=\frac12+\frac{27}{2}r_1r_2r_3s_1s_2s_3,\qquad
 \Phi=\frac14\max_{\substack{\varepsilon,\eta\in\{\pm1\}^3\\
                         \prod_j\varepsilon_j\eta_j=-1}}
 \left(1+\frac{\varepsilon\cdot r}{\sqrt3}\right)
 \left(1+\frac{\eta\cdot s}{\sqrt3}\right).
\end{equation}
Indeed, the four probabilities $|v_j|^2$ are
$(1+\sum_\alpha\chi_{j\alpha}r_\alpha s_\alpha)/4$, where
\[
 (\chi_{j\alpha})=
 \begin{pmatrix}1&-1&1\\1&1&-1\\-1&-1&-1\\-1&1&1\end{pmatrix}.
\]
Their second and third power sums in
\eqref{eq:fqe-global-intrinsic}, with $Z=D_v$ and $\operatorname{tr}(Z^TZ)=0$, give
the first identity. The 32 HS representatives become exactly the
auxiliary product rays whose Bloch vectors are cube vertices
$\varepsilon/\sqrt3,\eta/\sqrt3$ of opposite parity. The product-state
overlap formula and \eqref{eq:fqe-signed-diagonal-overlap} give the second
identity. These finite identities are checked independently by NC-1--NC-3.

Here is an elementary proof of the geometric claims in
\eqref{eq:fqe-null-critical-sharp}; NC-4 checks its scalar identities and constants.
Write $A=1-1/\sqrt3$, $B=1+1/\sqrt3$, and
$f_0=B^2/4=(2+\sqrt3)/6$.
First we prove $\Phi\ge f_0$. If the signs of the coordinates of
$r,s$ have opposite parity, align both cube vertices with those signs;
$\|r\|_1,\|s\|_1\ge1$ gives the claim. A zero coordinate allows either
parity. In the remaining case put $L=\|r\|_1$, $m=\min_j|r_j|$ and
define $f_r=1+L/\sqrt3$, $g_r=1+(L-2m)/\sqrt3$, with analogous
quantities for $s$. Flipping the sign at a smallest coordinate gives
\[
 4\Phi\ge\max(f_rg_s,g_rf_s)
                  \ge\sqrt{f_rg_rf_sg_s}.
\]
For the ordered coordinate moduli $x\ge y\ge z=m$ of either unit vector,
\[
 f_rg_r-B^2=\frac2{\sqrt3}(x+y-1)+\frac23(xy-z^2)\ge0.
\]
Here $xy\ge z^2$ and $(x+y)^2=1-z^2+2xy\ge1$.
This proves the fidelity floor in all cases.

If $r_1r_2r_3s_1s_2s_3\ge0$, then $Q\ge1/2$, and with $\Phi\ge f_0$ this gives
$Q\ge(1-\Phi)/(2(1-f_0))=\frac3{4-\sqrt3}(1-\Phi)$.
Otherwise put $\mu=3\sqrt3|r_1r_2r_3|$ and
$\nu=3\sqrt3|s_1s_2s_3|$, so $0\le\mu,\nu\le1$ and $2Q=1-\mu\nu$.
For a nonnegative unit triple $x,y,z$,
\[
 x+y+z-1
 =\frac{2(xy+yz+zx)}{x+y+z+1}
 \ge(9-3\sqrt3)xyz.
\]
For positive coordinates this follows from
$1/x+1/y+1/z\ge3\sqrt3$ and $x+y+z\le\sqrt3$;
zero coordinates follow by continuity.
Thus $\|r\|_1\ge1+(\sqrt3-1)\mu$ and similarly for $s$.
Aligning opposite-parity signs in \eqref{eq:fqe-null-critical-bloch} gives
$4\Phi\ge \Lambda:=(B+A\mu)(B+A\nu)$. The identity
\[
 \Lambda-B^2-(4-B^2)\mu\nu=AB\bigl(\mu(1-\nu)+\nu(1-\mu)\bigr)\ge0
\]
gives $1-\Phi\le(1-f_0)(1-\mu\nu)=2(1-f_0)Q$, the infidelity bound. In both cases the distance bound follows
pointwise: $1-\Phi=\frac{1+\sqrt\Phi}2\,d^2\ge\frac{1+\sqrt{f_0}}2\,d^2$, and
$\frac3{4-\sqrt3}\cdot\frac{1+\sqrt{f_0}}2=\frac{3+\sqrt3}4$ because $\sqrt{f_0}=B/2$.
The critical cluster representative $v=(1,i,0,0)/\sqrt2$ has
$Q=1/2$ and $\sqrt{\Phi}=B/2$, attaining all three geometric
bounds.


\begin{lemma}[Real-density geometry and sharp critical covering profile; exact certificate, RP, FP]
\label{lem:fqe-critical-real-density}
Write a unit critical normal vector as $v=x+iy$, with $x,y\in\mathbb R^4$,
and put $A=xx^T+yy^T$. Then $A$ ranges over exactly the real positive
semidefinite matrices of trace one and rank at most two, and
\begin{equation}\label{eq:fqe-critical-density-measurements}
 \lambda_{ji}=e_i^TH_jAH_j^Te_i,\qquad
 \tau=2\operatorname{tr}A^2-1.
\end{equation}
Here $H_0,H_1,H_2$ are the three real orthogonal cut transforms displayed
above. Their bases are mutually unbiased. For their twelve unit vectors
$h_{ji}=H_j^Te_i$ and rank-one projectors $B_{ji}=h_{ji}h_{ji}^T$,
\begin{equation}\label{eq:fqe-critical-real-tomography}
 A=\sum_{j,i}\lambda_{ji}B_{ji}-\tfrac12I,\qquad
 \sum_{j,i}\lambda_{ji}^2=\operatorname{tr}A^2+\tfrac12.
\end{equation}
For each coordinate axis $e_k$ and each of the four lines
$u=(0,\pm1,\pm1,\pm1)/\sqrt3$ in its orthogonal coordinate three-space,
put $P=e_ke_k^T+uu^T$ and $A_P=P/2$. There are sixteen such planes.
With $F_{\rm B}(A,B)=(\operatorname{tr}\sqrt{\sqrt B A\sqrt B})^2$,
\begin{equation}\label{eq:fqe-critical-bures-orbit}
 \Phi(v)=\max_P F_{\rm B}(A,A_P),\qquad
 F_{\rm B}(A,A_P)=\frac12\left[\operatorname{tr}(PA)
                    +2\sqrt{\det\bigl((e_k,u)^TA(e_k,u)\bigr)}\right].
\end{equation}
Thus $d(v)$ is exactly the minimum Bures distance from $A$ to these
sixteen real density matrices. The entire critical locus obeys the sharp bound
\begin{equation}\label{eq:fqe-critical-fidelity-floor}
 \Phi\ge f_{\rm cov}(\tau):=
       \frac13+\frac{\sqrt{1-\tau}}{2\sqrt3},\qquad
 \max_{\substack{\psi\ {\rm critical}\\\tau(\psi)=t}}d(\psi)^2
       =2-2\sqrt{f_{\rm cov}(t)}\quad(0\le t\le1).
\end{equation}
The bound is attained for every $t$ by the cluster-to-GHZ family below.
In particular every critical state has $d^2\le2-2/\sqrt3$, with equality for GHZ.
For any one such plane, put
$\lambda^P_{ji}=\operatorname{tr}(B_{ji}A_P)$ and
$D_P=\sum_{j,i}\lambda_{ji}\ln(\lambda_{ji}/\lambda^P_{ji})$.
The entropy gap has the exact form
\begin{equation}\label{eq:fqe-critical-relative-entropy}
 3\ln2\,\delta=D_P-\tfrac32\ln3\,J_P,\qquad
 J_P=1-\operatorname{tr}(PA)
    =\|A-A_P\|_F^2-\tau/2.
\end{equation}
\end{lemma}

\begin{proof}
The three physical cut vectors are $H_jv$, so their squared moduli are
the stated real measurements. Conversely, for any real $A$ of the
specified rank, its spectral decomposition gives
$v=\sqrt{a_1}e_1+i\sqrt{a_2}e_2$; rank one is included.
The identity $|v^Tv|^2=2\operatorname{tr}A^2-1$ proves the spin-flip claim.
The twelve real lines are the coordinate lines and the half-sign lines
of the regular $24$-cell. Their tight frame identity, valid for every
real symmetric matrix $C$, is
\[
 \sum_{j,i}\operatorname{tr}(B_{ji}C)B_{ji}
                    =C+\tfrac12\operatorname{tr}(C)I.
\]
Taking its trace pairing with $C$ proves both tomography formulas (RP-1).

The two HS rays associated with $P$ are $(e_k\pm iu)/\sqrt2$.
Their squared overlaps with $v$ have sum $\operatorname{tr}(PA)$ and
half their difference has squared value
$\det((e_k,u)^TA(e_k,u))$. Maximising over this pair therefore gives
\eqref{eq:fqe-critical-bures-orbit}, by the two-dimensional formula
$(\operatorname{tr}\sqrt C)^2=\operatorname{tr}C+2\sqrt{\det C}$.
The signed-Birkhoff reduction already proved shows that maximising
these sixteen pairwise maxima is the full physical HS fidelity.

For $q_\eta=|\langle\eta,v\rangle|^2$, the finite HS frame satisfies
\[
 \sum_{\eta\in\mathcal H}q_\eta=8,\qquad
 \sum_{\eta\in\mathcal H}q_\eta^2=\tfrac89(4-\tau).
\]
Both follow by expanding and cancelling odd sign powers (RP-3).
The stronger sharp profile \eqref{eq:fqe-critical-fidelity-floor} follows
from the area-weighted calibration proved below; the moment identity
alone does not give that profile.

Each reference distribution has one entry $1/2$ in each cut and nine
entries $1/6$ in total. The sum of the three projectors selecting its
large entries is $3P/2$ (RP-3). Expanding $D_P$ then gives
\eqref{eq:fqe-critical-relative-entropy}. Finally
$\operatorname{tr}A_P^2=1/2$ and
$2\operatorname{tr}(AA_P)=\operatorname{tr}(AP)$ give the leakage formula.
All matrix and compression identities here are independently checked by RP;
the covering calibration uses FP.
\end{proof}

\paragraph{Area-weighted covering calibration.}
We prove the profile in \eqref{eq:fqe-critical-fidelity-floor} using an
explicit probability measure on the sixteen HS planes. Its only
polynomial positivity input is the following finite exact certificate.

\begin{lemma}[Four-point cubic inequality; exact rational certificate]
\label{lem:fqe-four-point-cubic}
If $X,Y\in\mathbb R^4$ satisfy
$\sum_iX_i=\sum_iY_i=0$ and $\sum_iX_i^2=\sum_iY_i^2=4$, then
\begin{equation}\label{eq:fqe-four-point-cubic}
 \sum_{i,j=1}^4|X_i-Y_j|^3\ge64.
\end{equation}
Equality holds at $X=Y=(-1,-1,1,1)$.
\end{lemma}

\begin{proof}
Put $X=Er$, $Y=Es$, where
\[
 E=\begin{pmatrix}1&1&1\\1&-1&-1\\-1&1&-1\\-1&-1&1\end{pmatrix},
 \qquad E^TE=4I,
\]
so that $\|r\|=\|s\|=1$. Every zero-sum four-vector is represented.
Sort $X$ and $Y$ separately and interleave their entries in increasing
order. There are $\binom84=70$ words with four letters of each kind.
Swapping $X,Y$ and simultaneously negating both vectors reduce these
to exactly twenty-three classes (FP-1). Ties may use either order.
For each word let $z_1,\ldots,z_8$ be its ordered entries and
$g_i=z_{i+1}-z_i\ge0$. The absolute values then give a cubic polynomial
$p=\sum_{ij}|X_i-Y_j|^3/16-4$.
For each of the twenty-three words the supplied rational data give
the identity
\begin{equation}\label{eq:fqe-four-point-cubic-sos}
 p=B_0^TG_0B_0+\sum_{i=1}^7g_i B_i^TG_iB_i
       +(\|r\|^2-1)\ell_r+(\|s\|^2-1)\ell_s.
\end{equation}
The $\ell$ are quadratic polynomials and all eight $G_i$ are positive
definite. FP-2 rebuilds all $210$ coefficients through degree four
in each class and proves positivity of all $184$ Grams using integer
arithmetic. Thus the right side is nonnegative on each ordering cone
and the two spheres, proving the full inequality. The bases and exact
matrix representation are specified in Section~\ref{sec:sup-verifiers}.
The stated equality follows by substitution.
\end{proof}

Let $U=(U_1,U_2)$ be a real orthonormal $4\times2$ frame with plane projector
$\Pi=UU^T$. For each HS plane $P$ use its orthonormal frame
$V_P=(e_k,u)$ and put $c_P=|\det(U^TV_P)|$. The sixteen planes obey
\begin{equation}\label{eq:fqe-area-weighted-fusion}
 \sum_Pc_P^2=\frac83,\qquad
 \sum_Pc_P^2P=\frac89(I+\Pi).
\end{equation}
Indeed, the bivector $\Delta=U_1U_2^T-U_2U_1^T$ satisfies $\Delta\Delta^T=\Delta^T\Delta=\Pi$, and
$\det(U^TV_P)=e_k^T\Delta u$. Fix $k$, put $g=\Delta^Te_k$ (so $g_k=0$,
$|g|^2=\Pi_{kk}$) and sum over the four lines $u$ for this $k$, using $u_j=\pm1/\sqrt3$ ($j\ne k$):
$\sum_u(g\cdot u)^2=\frac43|g|^2$ and
$\sum_u(g\cdot u)^2uu^T=\frac49|g|^2(I-e_ke_k^T)+\frac89(gg^T-\operatorname{diag}(g_j^2))$.
Summing over $k$ with $\sum_kgg^T=\Pi$ and $\sum_k\operatorname{diag}(g_j^2)=\operatorname{diag}(\Pi)$ gives
both identities. FP-3 also checks them without
assuming orthonormality first, using the homogeneous squared area and
its associated matrix. Therefore the explicit weights
$w_P=3c_P^2/8$ satisfy
\begin{equation}\label{eq:fqe-plane-calibration}
 \sum_Pw_P=1,\qquad \sum_Pw_PU^TPU=\frac23I_2,\qquad
 \sum_Pw_Pc_P\ge\frac1{\sqrt3}.
\end{equation}

To prove the last inequality, let $M$ be the magic-basis matrix defined before
\eqref{eq:fqe-null-critical-bloch}. Form the null critical vector
$z=(U_1+iU_2)/\sqrt2$. Its auxiliary product representation
$Mz=a\otimes b$ has unit Bloch vectors $r,s$, as in
\eqref{eq:fqe-null-critical-bloch}.
For the two HS rays $\eta_P=(e_k+iu)/\sqrt2$ and $\bar\eta_P$,
expansion gives
\[
 c_P=\bigl||\langle\eta_P,z\rangle|^2
                 -|\langle\bar\eta_P,z\rangle|^2\bigr|.
\]
In the cube-pair description of the HS rays given after \eqref{eq:fqe-null-critical-bloch}, conjugation reverses both cube vertices.
Indeed $M=\epsilon^{\otimes2}\bar M$, where
$\epsilon=\left(\begin{smallmatrix}0&1\\-1&0\end{smallmatrix}\right)$
(FP-3), and the qubit spin flip reverses the Bloch vector.
Choose the four even-parity cube vertices $e_i$, the rows of $E$.
The sixteen conjugate pairs are represented once by
$(\varepsilon,\eta)=(e_i,-e_j)$, and hence
\[
 c_P=\frac{|e_i\cdot r-e_j\cdot s|}{2\sqrt3}.
\]
The four numbers $X_i=e_i\cdot r$, and likewise $Y_j=e_j\cdot s$,
have sum zero and sum of squares four. Lemma~\ref{lem:fqe-four-point-cubic}
therefore gives
\[
 \sum_Pc_P^3=\frac{\sum_{ij}|X_i-Y_j|^3}{24\sqrt3}
       \ge\frac8{3\sqrt3},
\]
which proves the last calibration inequality.

Now write the real density matrix of any critical state as
$A=U\Lambda U^T$, with $\operatorname{tr}\Lambda=1$ and
$\det\Lambda=(1-\tau)/4$. Rank one is included by completing its
support to any plane. The exact Bures formula gives
\begin{align*}
 \Phi
 &\ge\sum_Pw_P F_{\rm B}(A,P/2)\\
 &=\frac12\operatorname{tr}\!\left(\Lambda\sum_Pw_PU^TPU\right)
                  +\sqrt{\det\Lambda}\sum_Pw_Pc_P\\
 &\ge\frac13+\frac{\sqrt{1-\tau}}{2\sqrt3}.
\end{align*}
This proves the profile for the entire critical locus.

Finally, for $a,b\ge0$, $a+b>0$, the critical normal family
\begin{equation}\label{eq:fqe-covering-equality-family}
 v=\frac{(a+ib,-a+ib,0,0)}{\sqrt{2(a^2+b^2)}},\qquad
 \tau=\frac{(a^2-b^2)^2}{(a^2+b^2)^2},\qquad
 \Phi=\frac13+\frac{ab}{\sqrt3(a^2+b^2)}
\end{equation}
attains equality. FP-4 checks every one of the $32$ ray overlaps against
this maximum. The parameters cover $0\le\tau\le1$, proving pointwise
optimality and the conditional covering radius. The endpoint $a=b$
is a cluster state; $b=0$ is GHZ up to LU.


This model expresses both the entropy and the optimised orbit distance exactly.
In particular, a bound on the classical relative entropy $D_P$ that
controls the explicit plane leakage $J_P$ can improve the entropy
constant without first improving a moment surrogate. Its domain is a
real rank-two density matrix; replacing it by all real density matrices
would admit $I/4$, whose entropy exceeds the critical-state maximum.


\subsection{A one-parameter family of sharp geometric inequalities on the slice $P=1$}\label{sec:sup-nullsharp}
\begin{theorem}[Sharper geometric family on $P=1$; MS]\label{thm:sup-null-family}
Let $\psi$ be a critical unit pure four-qubit state with $\tau=0$, equivalently $P=1$, and let $\Phi,d,Q$ be as in
Theorem~\ref{thm:fqe-critical-orbit-sharp}. Use the two auxiliary Bloch vectors of \eqref{eq:fqe-null-critical-bloch} to put
$X=3\sum r_j^4-1$ and $Y=3\sum s_j^4-1$. Then
\begin{equation}\label{eq:fqe-null-critical-family}
 \begin{gathered}
 \Phi\ge\frac{2+\sqrt3}{6},\qquad 0\le\theta\le\frac1{24},\\
 Q+\theta XY\ge(1+8\theta)\frac3{4-\sqrt3}(1-\Phi)
       \ge(1+8\theta)\frac{3+\sqrt3}{4}d^2.
 \end{gathered}
\end{equation}
For every $\theta$, both geometric coefficients and the fidelity floor are optimal, attained by cluster states.
\end{theorem}
The case $\theta=0$ is the inequality \eqref{eq:fqe-null-critical-sharp} of the paper. The endpoint $\theta=\tfrac1{24}$ is proved by
an elementary multiaffine vertex check in the opposite-parity sector and a univariate polynomial minorant in the same-parity sector.
\paragraph{Sharp mixed-moment geometry.}
With the auxiliary Bloch vectors $r,s$ of
\eqref{eq:fqe-null-critical-bloch}, put
$X=3\sum r_j^4-1$ and $Y=3\sum s_j^4-1$.
We prove the endpoint $\theta=1/24$ in
\eqref{eq:fqe-null-critical-family} by elementary polynomial inequalities in both parity sectors.
Interpolation with the already proved $\theta=0$ bound, which is the inequality
\eqref{eq:fqe-null-critical-sharp} of the paper, then covers
$0\le\theta\le1/24$.

First, every nonnegative unit triple $a,b,c$ obeys the strengthened norm
inequality
\begin{equation}\label{eq:fqe-strengthened-sphere-norm}
 L:=a+b+c\ge1+(\sqrt3-1)u+\tfrac14(2-X-2u),\quad
 u=3\sqrt3abc,\quad X=3(a^4+b^4+c^4)-1.
\end{equation}
To see this, put $B=9-9\sqrt3/2$ and use
$a^4+b^4+c^4=1-(L^2-1)^2/2+4Labc$.
The desired residual is
\[
 L-1-\tfrac38(L^2-1)^2+(3L-B)abc.
\]
For $1\le L\le\sqrt2$, its first three terms equal
$(L-1)[1-3(L-1)(L+1)^2/8]\ge0$, since
$(\sqrt2-1)(\sqrt2+1)^2=1+\sqrt2<8/3$; also $3L>B$.
For $\sqrt2\le L\le\sqrt3$, Schur gives
$abc\ge L(L^2-2)/9$. The resulting lower bound is
\[
 f(L)=-L^4/24+(\sqrt3/2-1)L^3+L^2/12+(3-\sqrt3)L-11/8.
\]
It is concave on that interval, with
$f(\sqrt2)=\sqrt2-11/8>0$ and $f(\sqrt3)=0$, proving the claim.
All these finite identities are checked by MS-1.

Suppose first that $\prod_jr_js_j\le0$.
Put $u=3\sqrt3|r_1r_2r_3|$, $v=3\sqrt3|s_1s_2s_3|$.
For a unit triple $r$, put $t_i=r_i^2$ and $q_r=\sum_{i<j}t_it_j$. The identity
\[
 q_r^2-3t_1t_2t_3
 =\tfrac12\bigl[(t_1t_2-t_2t_3)^2+(t_2t_3-t_3t_1)^2
                         +(t_3t_1-t_1t_2)^2\bigr]\ge0
\]
gives the elementary bound
\begin{equation}\label{eq:fqe-bloch-triple-product}
 0\le u\le3q_r=1-X/2,\qquad 0\le v\le1-Y/2.
\end{equation}
Write $p=1-u$, $q=1-v$, $X=2p(1-\alpha)$ and
$Y=2q(1-\beta)$, with $p,q,\alpha,\beta\in[0,1]$; a zero
$p$ or $q$ is handled by continuity.
Aligning the two cube signs is admissible in this parity sector, so
\eqref{eq:fqe-strengthened-sphere-norm} gives
\[
 4\Phi\ge f_rf_s,\qquad
 f_r=2-(1-1/\sqrt3)p+\frac{p\alpha}{2\sqrt3},\quad
 f_s=2-(1-1/\sqrt3)q+\frac{q\beta}{2\sqrt3}.
\]
Set $a_{24}=4/(4-\sqrt3)$. The polynomial
\[
 \frac{p+q-pq}{2}+\frac{pq(1-\alpha)(1-\beta)}6
                   -a_{24}(1-f_rf_s/4)
\]
is affine in each of its four variables separately.
At the sixteen cube vertices its distinct values are
\[
 0,\quad \frac{3+4\sqrt3}{26},\quad
 \frac{-6+5\sqrt3}{78},\quad
 \frac{2+7\sqrt3}{52},\quad \frac{3+4\sqrt3}{78},
\]
all nonnegative (MS-1). This proves
$Q+XY/24\ge a_{24}(1-\Phi)$ in this sector.

For the remaining parity sector ($\prod_{j=1}^3 r_j s_j\ge0$), sort the
coordinate moduli of $r$ as $a_r\ge b_r\ge c_r\ge0$ and those of $s$ as
$a_s\ge b_s\ge c_s\ge0$. Swap $r$ and $s$ if necessary so that
$c_r\ge c_s$, and set
\[
 x=\sqrt3\,c_r\in[0,1],\quad y=\sqrt3\,c_s\in[0,x],\qquad
 \tilde w_r=3a_rb_r,\quad \tilde w_s=3a_sb_s.
\]
For a sorted unit triple $a\ge b\ge c\ge0$ with $x=\sqrt3\,c$ and
$\tilde w=3ab$, one has $a+b=\sqrt{1-c^2+2ab}=\sqrt{1+(2\tilde w-x^2)/3}$,
$3\sqrt3\,abc=x\tilde w$, and
$\tilde w\in\bigl[x\sqrt{3-2x^2},\,(3-x^2)/2\bigr]$, where the lower and
upper bounds correspond to $b=c$ and $a=b$. Moreover (MS-2),
\[
 X=2-2x^2\Bigl(1-\frac{x^2}{3}\Bigr)-\frac23\tilde w^2
  \ge X_{\rm lin}(x,\tilde w)
  :=2-2x^2\Bigl(1-\frac{x^2}{3}\Bigr)-\Bigl(1-\frac{x^2}{3}\Bigr)\tilde w\ge0,
\]
since $X-X_{\rm lin}=\frac23\tilde w\bigl(\frac{3-x^2}{2}-\tilde w\bigr)=3ab(a-b)^2\ge0$
and $X_{\rm lin}(x,\tilde w)=\frac12(1-x^2)^2+\bigl(1-\frac{x^2}{3}\bigr)\bigl(\frac{3-x^2}{2}-\tilde w\bigr)\ge0$.
Using the branch $(1+(a_r+b_r+c_r)/\sqrt3)(1+(a_s+b_s-c_s)/\sqrt3)$ of $4\Phi$ gives
$4\Phi\ge f_r(x,\tilde w_r)\,g_s(y,\tilde w_s)$ with
\[
 f_r(x,\tilde w_r)=1+\frac1{\sqrt3}\sqrt{1+\frac{2\tilde w_r-x^2}{3}}+\frac x3,\qquad
 g_s(y,\tilde w_s)=1+\frac1{\sqrt3}\sqrt{1+\frac{2\tilde w_s-y^2}{3}}-\frac y3.
\]
Hence $Q+XY/24-a_{24}(1-\Phi)$ is bounded below by
\[
 \mathcal F(x,y,\tilde w_r,\tilde w_s)
 :=\frac{1+xy\tilde w_r\tilde w_s}{2}
   +\frac{X_{\rm lin}(x,\tilde w_r)\,Y_{\rm lin}(y,\tilde w_s)}{24}
   +\frac{a_{24}}{4}f_r(x,\tilde w_r)\,g_s(y,\tilde w_s)-a_{24}.
\]
Since $(3-2x^2)-\bigl(\frac{5-2x}{3}\bigr)^2=\frac{2(1-x)(1+11x)}{9}\ge0$,
we have $\tilde w_r\ge x\sqrt{3-2x^2}\ge\tilde w_{\rm low}(x):=\frac{x(5-2x)}{3}\ge x$
and $\tilde w_s\ge\tilde w_{\rm low}(y)\ge y$. Consequently $2\tilde w_s-y^2\ge\frac{y(10-7y)}{3}\ge0$, so
\[
 Y_{\rm lin}(y,\tilde w_s)\le2-\frac23y
 \le\frac{2\sqrt3}{\sqrt3+1}\Bigl(1+\frac1{\sqrt3}-\frac y3\Bigr)
 \le\frac{2\sqrt3}{\sqrt3+1}\,g_s(y,\tilde w_s),
\]
and likewise $X_{\rm lin}(x,\tilde w_r)\le2\le\frac{2\sqrt3}{\sqrt3+1}\,f_r(x,\tilde w_r)$.
Since $1+(2\tilde w_r-x^2)/3\le2$, the partial derivative with respect to $\tilde w_r$ satisfies
\[
 \frac{\partial\mathcal F}{\partial\tilde w_r}
 \ge Y_{\rm lin}(y,\tilde w_s)\Bigl(\frac{\sqrt3+1}{18\sqrt2\,(4-\sqrt3)}-\frac1{24}\Bigr)\ge0,
\]
and the same calculation gives $\partial\mathcal F/\partial\tilde w_s\ge0$ (MS-2). Monotonicity in $(\tilde w_r,\tilde w_s)$
therefore reduces the inequality to $\tilde w_r=\tilde w_{\rm low}(x)$ and
$\tilde w_s=\tilde w_{\rm low}(y)$.

At $\tilde w=\tilde w_{\rm low}(x)$ one has $x\tilde w_{\rm low}(x)=\frac{x^2(5-2x)}{3}\ge x^2$,
\[
 X_{\rm lin}\bigl(x,\tilde w_{\rm low}(x)\bigr)-\frac49x^2(1-x)^2
 =\frac19(1-x)(18+3x-13x^2)=:X_b(x)\ge0,
\]
and $\sqrt{1+(2\tilde w_{\rm low}(x)-x^2)/3}=\frac13\sqrt{9+10x-7x^2}\ge S_b(x):=1+\frac12x-\frac7{20}x^2$,
since $(9+10x-7x^2)-9S_b(x)^2=\frac{x(10-7x)(63x^2-90x+40)}{400}\ge0$ on $[0,1]$ (MS-2).
Putting $\hat f(x)=1+S_b(x)/\sqrt3+x/3$ and $\hat g(y)=1+S_b(y)/\sqrt3-y/3$, we obtain
\begin{equation}\label{eq:fqe-null-mixed-certificate}
 Q+\frac{XY}{24}-a_{24}(1-\Phi)
 \ge P(x,y):=\frac{1+x^2y^2}{2}+\frac{X_b(x)X_b(y)}{24}+\frac{a_{24}}{4}\hat f(x)\hat g(y)-a_{24}
\end{equation}
for $0\le y\le x\le1$. Write $B=1+1/\sqrt3$. Three univariate polynomials on $[0,1]$ close the bound (MS-3):
\begin{enumerate}
\item $\hat g(y)-\frac B2 X_b(y)=\frac y2 q_2(y)\ge0$, where
 $q_2(y)=\bigl(1+\frac{8\sqrt3}{9}\bigr)+\bigl(\frac{16}{9}+\frac{97\sqrt3}{270}\bigr)y-\bigl(\frac{13}{9}+\frac{13\sqrt3}{27}\bigr)y^2$
 is concave on $[0,1]$ with $q_2(0)>0$ and $q_2(1)=\frac43+\frac{23\sqrt3}{30}>0$.
\item Consequently, for $0\le y\le x\le1$,
 \[
  \frac{\partial P}{\partial x}(x,y)
  =xy^2+\frac{a_{24}}{4}\hat f'(x)\hat g(y)+\frac{X_b'(x)}{24}X_b(y)
  \ge\frac{X_b(y)}{4-\sqrt3}\,q_1(x)\ge0,
 \]
 where $q_1(x):=(4-\sqrt3)\bigl(\frac{a_{24}B}{8}\hat f'(x)+\frac{X_b'(x)}{24}\bigr)
 =\bigl(\frac{11}{20}-\frac35x\bigr)^2+\bigl(\frac{5\sqrt3}{24}-\frac{1189}{3600}\bigr)
 +\bigl(\frac{17\sqrt3}{540}-\frac{133}{2700}\bigr)x+\bigl(\frac{163}{450}-\frac{13\sqrt3}{72}\bigr)x^2>0$ on $[0,1]$.
\item Hence $P(x,y)\ge P(y,y)=yq_3(y)\ge0$ on $[0,1]$, where
 \begin{align*}
  (4-\sqrt3)\,q_3(y)
  &=\Bigl(\frac{13}{25}-\frac{26}{25}y-\frac1{10}y^2\Bigr)^2
   +\Bigl(\frac{11\sqrt3}{18}-\frac{5896}{5625}\Bigr)
   +\Bigl(\frac{737}{7500}-\frac{19\sqrt3}{360}\Bigr)y\\
  &\quad+\Bigl(\frac{173411}{202500}-\frac{79\sqrt3}{162}\Bigr)y^2
   +\Bigl(\frac{756757}{486000}-\frac{419\sqrt3}{972}\Bigr)y^3\\
  &\quad-\Bigl(\frac{21043}{24300}-\frac{52\sqrt3}{243}\Bigr)y^4
   +\frac{169(4-\sqrt3)}{1944}y^5>0,
 \end{align*}
 since every coefficient through $y^3$ and at $y^5$ is strictly positive while the $y^3$ coefficient ($>4/5$)
 exceeds the $y^4$ coefficient ($<1/2$) on $0\le y\le1$. Since $P(0,0)=0$, this proves $P(x,y)\ge0$ on the entire triangle $0\le y\le x\le1$.
\end{enumerate}

Finally,
\[
 Q+\theta XY=(1-24\theta)Q+24\theta(Q+XY/24)
 \ge(1+8\theta)\frac3{4-\sqrt3}(1-\Phi).
\]
The fidelity floor already proved converts this to the displayed distance
bound. Clusters have $Q=1/2$, $X=Y=2$ and
$\Phi=(2+\sqrt3)/6$, attaining both coefficients for every
$\theta$ in the stated interval (MS-4). Thus the full geometric family
is sharp.


\subsection{The minimum-purity slice}
\label{sec:fqe-p1-sharp}
Statements are proved; the finite verifications are interval-certified (outward-rounded ball
arithmetic, every sign decision a certain comparison) or exact, and are listed in
Section~\ref{sec:sup-verifiers} (PM, PS).
This section gives the coordinates, symmetries and equality geometry on the slice $P=1$ ($\tau=0$) that
Section~\ref{sec:fqe-critical-sharp} uses for the whole critical locus.

\paragraph{Geometry of the equality states.}
At an equality point $\delta$ is stationary in every direction, because one Bloch vector is an axis
 and $\delta$ is even in the vanishing coordinate of the other, while four
admissible cube pairs tie (PM-4) and, computationally, their $d^2$-gradients form a regular tetrahedron around $0$ (explanatory; not used). The distance $d^2$
therefore has a conical corner and $\delta-C_*d^2$ has a conical minimum in the three transverse directions;
the one smooth direction, which rotates the second vector within its coordinate plane, is the wall of
Lemma~\ref{lem:fqe-p1-wall}, where the wall function $G_w$ is stationary.  This picture is explanatory; the proof uses the
exact identities of Lemma~\ref{lem:fqe-critical-local-cone}.

\subsubsection*{Coordinates and inputs}
Let $r,s$ be the unit auxiliary Bloch vectors of \eqref{eq:fqe-null-critical-bloch}; indices are taken mod $3$.
\begin{lemma}[Cut spectra; PM-1]\label{lem:fqe-p1-cuts}
For $k\in\mathbb Z/3$ put $u^{(k)}_a=r_{a+k}s_a$ and let
\[
 \mathcal S=\{(-,-,-),(-,+,+),(+,-,+),(+,+,-)\}.
\]
The three physical cuts have the spectra $\{(1+\sigma\cdot u^{(k)})/4:\sigma\in\mathcal S\}$, $k=0,1,2$ in some order.
Hence $E=\frac13\sum_kg(u^{(k)})$, $g(u)=\sum_{\sigma\in\mathcal S}\eta\bigl((1+\sigma\cdot u)/4\bigr)$,
$\eta(x)=-x\log_2x$, $\delta=1+\tfrac12\log_23-E$, and $\sum_k|u^{(k)}|^2=|r|^2|s|^2=1$.
\end{lemma}
\begin{proof}
Let $h_{ji}=H_j^Te_i$ be the unit vectors of Lemma~\ref{lem:fqe-critical-real-density}. Since $h_{ji}$ is real,
$Mh_{ji}$ is a Bell ray; write
$C^{ji}_{\alpha\beta}=\langle Mh_{ji}|\sigma_\alpha\otimes\sigma_\beta|Mh_{ji}\rangle$ for its correlation matrix, which is real
orthogonal with determinant $-1$. Its projector is $\tfrac14(I+\sum_{\alpha\beta}C^{ji}_{\alpha\beta}\sigma_\alpha\otimes\sigma_\beta)$, so
the probability $\lambda_{ji}=|h_{ji}^Tv|^2=|\langle Mh_{ji},Mv\rangle|^2$ is $(1+r^TC^{ji}s)/4$, as $Mv=a\otimes b$.
\emph{Cut $0$.} From the definitions of the $m_k$, $(\sigma_\alpha\otimes\sigma_\alpha)m_k=s^\alpha_km_k$ with
$s^x=(+,+,-,-)$, $s^y=(-,+,-,+)$, $s^z=(+,-,-,+)$. So $C^{0i}$ has diagonal $(s^x_i,s^y_i,s^z_i)$, entries $\pm1$, and being
orthogonal it is diagonal; the four diagonals $(+,-,+),(+,+,-),(-,-,-),(-,+,+)$ are the elements of $\mathcal S$.
\emph{Cuts $1,2$.} Let $u=\tfrac12\bigl(-I+i(\sigma_z-\sigma_x-\sigma_y)\bigr)$, a unitary with $u\sigma_\beta u^\dagger=\sum_\alpha
R_{\alpha\beta}\sigma_\alpha$, where $Re_1=e_2$, $Re_2=-e_3$, $Re_3=-e_1$. Direct multiplication gives
$MH_2M^\dagger=u\otimes u^\dagger$ and $MH_1M^\dagger=(u\otimes u^\dagger)\,\mathrm{SWAP}$, where SWAP exchanges the two auxiliary
qubits and transposes correlation matrices. Since $Mh_{ji}=(MH_jM^\dagger)^\dagger m_i$, this gives
$C^{2i}=R^TC^{0i}R^T$ and $C^{1i}=RC^{0i}R$. As $R_{\alpha\beta}\ne0$ only for $\alpha=\beta+1$, these are the signed
permutations $C^{ji}_{\alpha\beta}=\sigma_\beta[\alpha=\beta+k_j]$ with $(k_j)=(0,2,1)$; the four column-sign triples $\sigma$ of
each cut are distinct and have product $\det C^{ji}=-1$, so they run over $\mathcal S$. Hence
$r^TC^{ji}s=\sigma\cdot u^{(k_j)}$. PM-1 recomputes all twelve $C^{ji}$ exactly from the physical magic basis.
\end{proof}
\begin{lemma}[Symmetries; PM-2]\label{lem:fqe-p1-sym}
$\delta$ and $d^2$ are invariant under: even sign changes of the six coordinates; cyclic shifts of $r$ alone and of $s$
alone; and the exchange $r\leftrightarrow s$.  (The stated operations permute both the admissible cube pairs and the twelve
probabilities. Fidelity alone is also invariant under arbitrary coordinate permutations
of either vector.)
\end{lemma}

\subsubsection*{The wall}
Let $h(t)=H_{\rm b}((1+t)/2)$, with $H_{\rm b}$ the binary entropy defined in Appendix~\ref{sec:fqe-critical-orbit-stability}.
For $r=e_1$ and $s=(\cos\theta,0,\sin\theta)$ (a point of the wall) write
$\delta_w(\theta)=\frac12\log_23-\frac13(1+h(\cos\theta)+h(\sin\theta))$,
$d_w^2(\theta)=2-\sqrt{B(1+(\cos\theta+\sin\theta)/\sqrt3)}$, $B=1+1/\sqrt3$, and $G_w=\delta_w-C_*d_w^2$.
\begin{lemma}[Wall]\label{lem:fqe-p1-wall}
$G_w(\pi/4)=0$ by \eqref{eq:fqe-p1-cstar}, and $G_w(\theta)=G_w(\pi/2-\theta)$; hence $G_w'(\pi/4)=0$.
\end{lemma}
\begin{proof}
The reflection $\theta\mapsto\pi/2-\theta$ exchanges $\cos\theta$ and $\sin\theta$, so it exchanges $h(\cos\theta)$ and
$h(\sin\theta)$ and fixes $\cos\theta+\sin\theta$; thus $\delta_w$ and $d_w^2$ are both invariant.  $G_w$ is smooth near
$\pi/4$ (there $\cos\theta,\sin\theta\in(0,1)$, where $h$ is analytic), and an even function of $\theta-\pi/4$ has zero
derivative at $\pi/4$.
\end{proof}


\subsection{The sharp entropy coefficient on the whole critical locus}
\label{sec:fqe-critical-sharp}
Statements are proved; the finite verifications are interval-certified (outward-rounded ball arithmetic, every sign
decision a certain comparison) or exact, and are listed in Section~\ref{sec:sup-verifiers} (TC, EX, LCC, CCV).
This section proves Theorem~\ref{thm:fqe-p1-sharp} for all critical states.
Section~\ref{sec:fqe-p1-sharp} supplies the coordinates, symmetries, wall function and equality geometry on the slice
$\tau=0$ ($P=1$) that are used here.

The constant $C_*$ is defined in \eqref{eq:fqe-p1-cstar}; the equality states attain it exactly (PM-4), and PS-1 encloses its decimal value.

\begin{restated}[Sharp entropy stability on the critical locus; interval-certified]{Theorem}{thm:fqe-p1-sharp}
Every critical unit pure four-qubit state satisfies
\[
 \delta(\psi)\ \ge\ C_*\,d(\psi)^2 ,
\]
and equality with $d>0$ holds exactly for states whose critical normal form has
one of the following $144$ Bloch-coordinate pairs, all with $P=1$. In the auxiliary Bloch coordinates $(r,s)\in S^2\times S^2$ of the null-critical normal form they are: one of $r,s$ is a signed
coordinate axis and the other is $(\pm e_b\pm e_c)/\sqrt2$, $b\ne c$ ($6\cdot12$ pairs, twice for the swap).
These pairs attain the bound, so $C_*$ is the optimal coefficient on the critical locus. Consequently every universal coefficient in
$\delta\ge Cd^2$ is at most $C_*$, and the question left open is whether $C_*$ is optimal over all states.
\end{restated}

\paragraph{Strategy.}
On the critical locus the twelve cut probabilities are linear in a $3\times3$ correlation matrix $T$, and
$\delta=\delta(T)$ is convex. The fidelity is the maximum of sixteen explicit plane functions. A critical state is
determined by $(r,s,w)$ with unit vectors $r,s$ (the $P=1$ coordinates) and a disc coordinate $w\in\mathbb R^2$,
$|w|^2=\tau$; the equality states sit at $w=0$. Analytic bounds settle $\tau\ge0.36$ (their crossing is at $\tau\approx0.35879$, TC-6); for $\tau\le0.37$ a
six-dimensional interval covering, together with a local cone lemma at the equality points, closes the proof.

\subsubsection*{Critical states as correlation matrices}
Let $M$ be the magic matrix, $v$ a unit critical normal vector, $A=\operatorname{Re}(vv^\dagger)$, and
$\psi=Mv$ the auxiliary two-qubit vector.  Put $T_{\alpha\beta}=\langle\psi|\sigma_\alpha\otimes\sigma_\beta|\psi\rangle$,
$t_{ji}=4\lambda_{ji}-1$ with $\lambda_{ji}$ the cut probabilities of Lemma~\ref{lem:fqe-critical-real-density}.
\begin{lemma}[Correlation-matrix model; analytic, exact check TC-1]\label{lem:fqe-critical-correlation}
\begin{enumerate}
\item The twelve rays $Mh_{ji}$ are Bell rays without local terms; their correlation matrices $C^{ji}$ are signed
permutation matrices (those of Lemma~\ref{lem:fqe-p1-cuts}), and $t_{ji}=\langle C^{ji},T\rangle$.
Consequently, by \eqref{eq:fqe-critical-real-tomography},
$A=\tfrac14I+\tfrac14\sum_{j,i}t_{ji}B_{ji}$: $A$ is an affine function of $T$.
\item $T=R_a\operatorname{diag}(\kappa,-\kappa,1)R_b^{T}$ with $R_a,R_b\in SO(3)$ and $\kappa=\sqrt\tau$.  Writing
$r=R_ae_3$, $s=R_be_3$ and $x_a,y_a$ ($x_b,y_b$) for the first two columns of $R_a$ ($R_b$),
\[
 T=rs^T+\kappa\bigl(x_ax_b^T-y_ay_b^T\bigr).
\]
Rotating the frames about $r$ and $s$ by $\theta_a,\theta_b$ turns $\kappa(x_ax_b^T-y_ay_b^T)$ into
$w_1(x_ax_b^T-y_ay_b^T)+w_2(x_ay_b^T+y_ax_b^T)$ with $w=\kappa(\cos(\theta_a+\theta_b),\sin(\theta_a+\theta_b))$.
\item For a plane $\Pi=VV^T$ of Lemma~\ref{lem:fqe-critical-real-density} the matrix $m_\Pi=V^TAV$ is affine in $T$ and
$\Phi=\max_\Pi\Phi_\Pi$, $\Phi_\Pi=\tfrac12\operatorname{tr}m_\Pi+\sqrt{\det m_\Pi}$.  Equivalently
$\Phi=\max\tfrac14\bigl(1+\sqrt{1-\tau}\,(\varepsilon\!\cdot\! r+\eta\!\cdot\! s)/\sqrt3+\varepsilon^TT\eta/3\bigr)$ over the
$32$ cube pairs of \eqref{eq:fqe-null-critical-bloch}.
\end{enumerate}
\end{lemma}
\begin{proof}
(1) The projector onto a maximally entangled ray is $\tfrac14(I+\sum C_{\alpha\beta}\sigma_\alpha\otimes\sigma_\beta)$,
with no local terms; the proof of Lemma~\ref{lem:fqe-p1-cuts} derives the twelve $C^{ji}$, and $t_{ji}=\langle C^{ji},T\rangle$
because $\lambda_{ji}$ is the expectation of that projector in $\psi$ (PM-1 and TC-1 recompute them, TC-1 together
with the identity $\lambda_{ji}=(\|v\|^2+\langle C^{ji},T\rangle)/4$ exactly for a random rational $v$).
(2) A pure two-qubit vector is $(U\otimes V)(\cos\varphi|00\rangle+\sin\varphi|11\rangle)$ with $0\le\varphi\le\pi/4$; its
correlation matrix is $\operatorname{diag}(\sin2\varphi,-\sin2\varphi,1)$ in the frames rotated by $U,V$, and the concurrence
$\sin2\varphi=|\sum_jv_j^2|=\sqrt\tau$ (the coefficient matrix of $Mv$ has determinant $\sum_jv_j^2/2$).  The local Bloch
vectors are $\cos2\varphi\,r$ and $\cos2\varphi\,s$, of length $\sqrt{1-\tau}$.
(3) $m_\Pi$ is affine in $A$.  The first formula is \eqref{eq:fqe-critical-bures-orbit}.  The second is the overlap of
$\psi$ with a product ray with Bloch vectors $\varepsilon/\sqrt3,\eta/\sqrt3$:
$\langle\psi|\tfrac12(1+\hat\varepsilon\cdot\sigma)\otimes\tfrac12(1+\hat\eta\cdot\sigma)|\psi\rangle
=\tfrac14(1+\hat\varepsilon\cdot r_A+\hat\eta\cdot r_B+\hat\varepsilon^TT\hat\eta)$.
TC-2 and TC-3 check the affine representation of $A$, the equality of the two maxima for 200 random states, and the singular
values $(1,\kappa,\kappa)$, $\det T=-\kappa^2$.
\end{proof}
Thus $\tau=|w|^2$, and a critical state is determined (up to the discrete choices) by $(r,s,w)$.  At $\tau=0$ one recovers the
auxiliary Bloch vectors of \eqref{eq:fqe-null-critical-bloch}.  In these variables
\[
 \begin{gathered}
 \delta=\mathrm{EMAX}-\tfrac13\sum_{j,i}\operatorname{ent}\bigl(\tfrac{1+t_{ji}}4\bigr),\qquad
 d^2=2-2\sqrt{\max_\Pi\Phi_\Pi},\\
 G:=\delta-C_*d^2=\max_\Pi G_\Pi,\qquad G_\Pi=\delta-C_*(2-2\sqrt{\Phi_\Pi}),
 \end{gathered}
\]
with $\mathrm{EMAX}=1+\frac12\log_23$ and $\operatorname{ent}(x)=-x\log_2x$.  Each $G_\Pi$ is analytic wherever all $\lambda_{ji}>0$ and
$\det m_\Pi>0$, and $\delta$ is a convex function of $T$.

\subsubsection*{The chart}
For $\zeta_1,\zeta_2,\zeta_1',\zeta_2'\in[0,1]$ and a sign $\mathrm{sg}=\pm1$ put
\begin{align*}
 r&=\tfrac1N(1,\zeta_1,\zeta_2),\ x_a=\tfrac1n(\zeta_1,-1,0),\ y_a=r\times x_a,\\
 s&=\tfrac1{N'}(\mathrm{sg},\zeta_1',\zeta_2'),\ x_b=\tfrac1{n'}(\zeta_1',-\mathrm{sg},0),\ y_b=s\times x_b,
\end{align*}
$N=\sqrt{1+\zeta_1^2+\zeta_2^2}$, $n=\sqrt{1+\zeta_1^2}$ (and primed), and for $w\in\mathbb R^2$, $|w|\le1$,
\begin{equation}\label{eq:fqe-critical-chart}
 T(y)=rs^T+w_1(x_ax_b^T-y_ay_b^T)+w_2(x_ay_b^T+y_ax_b^T),\qquad y=(\zeta_1,\zeta_2,\zeta_1',\zeta_2',w_1,w_2).
\end{equation}
\begin{lemma}[Reduction to the chart; analytic]\label{lem:fqe-critical-chart}
$\delta$ and $d$ are invariant under $T\mapsto S_aTS_b^T$, where $S_a=D_a\Pi_a$, $S_b=D_b\Pi_b$ with
$\Pi_a,\Pi_b$ cyclic permutation matrices and $D_a,D_b$ diagonal sign matrices, $\det D_a\det D_b=1$ (an even total number of
sign changes), and under the transposition $T\mapsto T^T$.  Every critical state is equivalent under such maps to one
with $T=T(y)$, $y$ in the chart $[0,1]^4\times\{|w|\le1\}$ and $\mathrm{sg}=\pm1$, and then $\tau=|w|^2$.
In the chart the transposition is $(\zeta_1,\zeta_2,\zeta_1',\zeta_2',w_1,w_2)\mapsto(\zeta_1',\zeta_2',\zeta_1,\zeta_2,w_1,\mathrm{sg}\,w_2)$, so it suffices to treat $\zeta_1\le\zeta_1'$.
\end{lemma}
\begin{proof}
The $t_{ji}$ are permuted: this is a linear identity in $T$, and holds on the rank-one matrices $rs^T$ by
Lemma~\ref{lem:fqe-p1-sym}, which span all matrices.  The set of $32$ cube pairs is invariant (PM-2: the parity
$\prod\varepsilon_j\eta_j$ is preserved because $\det D_a\det D_b=1$), the local vectors transform with $S_a,S_b$ up to a
common sign which is absorbed by $(\varepsilon,\eta)\mapsto-(\varepsilon,\eta)$, and $\varepsilon^TT\eta$ is invariant;
so $\Phi$ is invariant by Lemma~\ref{lem:fqe-critical-correlation}(3).  Given $T$, take $r,s$ as in
Lemma~\ref{lem:fqe-critical-correlation}(2); cyclic shifts of $r$ alone and of $s$ alone make $|r_1|,|s_1|$ maximal, and
flips of an even number of the six coordinate signs make $r_1,r_2,r_3,s_2,s_3\ge0$ (the sign of $s_1$ is then determined and
is $\mathrm{sg}$).  The sign patterns of $r$ and of $s$ may have $\det D_a=\det D_b=-1$; then $D_a=-D_a'$, $D_b=-D_b'$ with
$\det D_a'=\det D_b'=1$ and $D_aTD_b^T=D_a'TD_b'^T$, i.e.\ the common sign flip of $(r,s)$, which leaves $rs^T$ fixed, is absorbed.
$(r,s)$ is determined by $T$ up to this common sign when $\tau<1$; at $\tau=1$ the top singular value is degenerate and $r,s$ are
not unique, but $\tau=1$ lies in the analytic range and the chart is used only for $\tau\le0.37$.  Then $r=(1,\zeta_1,\zeta_2)/N$, $s=(\mathrm{sg},\zeta_1',\zeta_2')/N'$ with $\zeta_1,\zeta_2,\zeta_1',\zeta_2'\in[0,1]$, and $T-rs^T$ lies in the
span of the two matrices of \eqref{eq:fqe-critical-chart} by Lemma~\ref{lem:fqe-critical-correlation}(2), with $|w|=\kappa$.
Transposition permutes the twelve $t_{ji}$ (again a linear identity, on rank-one matrices Lemma~\ref{lem:fqe-p1-sym} gives the swap
$r\leftrightarrow s$) and exchanges the roles of the two cube vertices in $\Phi$. Its action on the chart coordinates follows by
direct substitution: $T(y)^T=DT(y')D$, with $y'$ the image coordinates displayed in the lemma and $D=I$ for $\mathrm{sg}=1$,
$D=\operatorname{diag}(1,-1,-1)$ for $\mathrm{sg}=-1$, an admissible sign change (TC-4 also checks it numerically).
\end{proof}

\subsubsection*{Large $\tau$}
\begin{lemma}[Analytic range; TC-6]\label{lem:fqe-critical-large-tau}
If $\tau\ge0.36$ then $\delta\ge C_*d^2$, strictly.
\end{lemma}
\begin{proof}
Put $k=(\ln3-1)/(2\ln2)$.  By \eqref{eq:fqe-critical-entropy} and \eqref{eq:fqe-critical-fidelity},
$\delta\ge k(1-\Phi+5\tau/2)$, and $\Phi\ge f_{\rm cov}(\tau)$ by \eqref{eq:fqe-critical-fidelity-floor}.  With
$x=\sqrt{\Phi}$ the function $k(1-x^2)+5k\tau/2-2C_*(1-x)$ is concave in $x$, hence on $[\sqrt{f_{\rm cov}},1]$ it
is at least its smaller endpoint value; at $x=1$ it is $5k\tau/2>0$.  At $x=\sqrt{f_{\rm cov}(\tau)}$, writing
$\tau=1-u^2$, $f=f_{\rm cov}=\frac13+\frac u{2\sqrt3}$, it is
$\psi(u)=k(1-f)+\frac52k(1-u^2)-2C_*(1-\sqrt f)$, which is concave on $0\le u\le4/5$: $k(1-f)$ is affine,
$-\frac52ku^2$ is concave, and $\sqrt f$ is the square root of a positive affine function, hence concave, with coefficient
$2C_*>0$.  So $\psi$ is positive on $[0,4/5]$ once it is positive at the two endpoints, which TC-6 encloses in ball
arithmetic ($\psi(0)>0.064$, $\psi(4/5)>1.7\cdot10^{-4}$).
\end{proof}

\subsubsection*{The local cone lemma}
At $w=0$ the equality points of Theorem~\ref{thm:fqe-p1-sharp} lie in the chart at the eight points
$(\zeta_1,\zeta_2,\zeta_1',\zeta_2')\in\{(0,0,0,1),(0,0,1,0),(0,1,0,0),(1,0,0,0)\}$, one set for each sign $\mathrm{sg}$.
Let $z_i=|y_i-y^{\rm eq}_i|$ for the four chart coordinates and $z_{5,6}=w_{1,2}$; the box $\mathcal B_\rho$ is
$|z|_\infty\le\rho$, $\rho=1/256$.
\begin{lemma}[Local cone; LCC]\label{lem:fqe-critical-local-cone}
At each equality point exactly four planes $\Pi$ are tied: $G_\Pi=0$ and the others are below $-0.05$.  Their gradients (in $z$)
sum to $0$ and have rank three.  In $\mathcal B_\rho$ one has $G>0$ except at the point itself.
\end{lemma}
\begin{proof}
By Taylor's formula with the Hessian enclosed over the box, $G_\Pi(t\omega)=t\,g_\Pi\!\cdot\!\omega+\frac{t^2}2\omega^T H_\Pi(\xi)\omega$
for $\omega$ on the unit sphere of $|\cdot|_\infty$ (in the orthant of the four chart coordinates) and $0<t\le\rho$, with
$H_\Pi(\xi)$ in an interval matrix $[H_\Pi]$, and $G\ge\sum_\Pi\mu_\Pi G_\Pi$ for weights $\mu_\Pi\ge0$ summing to one on the tied planes.
Write $\mu=\frac14+\nu$.  Because $\sum g_\Pi=0$, the first-order term of the average depends on $\nu$ only.  The tie and this identity are exact, and
are proved symbolically (EX, no floating point or enclosures) at all eight points: exactly four planes have $\Phi_\Pi=\Phi_*$, where
$4\Phi_*=(1+1/\sqrt3)(1+\sqrt{2/3})$ and all other $\Phi_\Pi$ are smaller; the twelve probabilities are $(2\pm\sqrt2)/8$ and $\frac14$, so
$\delta=\frac12\log_23-\frac13(1+2H_{\rm b}(p))$ and $\delta=C_*(2-2\sqrt{\Phi_*})$, i.e.\ $G_\Pi=0$ at the point for the four tied planes, which
is the definition of $C_*$ in \eqref{eq:fqe-p1-cstar}; and $\sum_{\Pi\ \rm tied}\partial_kG_\Pi=0$ for all six $k$, as an identity in
$C_*,\ln p,\ln(1-p),\ln2$.  The reason is the sign symmetry
$z\mapsto(s_1z_1,s_2z_2,s_1s_2z_3,z_4,s_2z_5,s_1z_6)$ at $(0,0,0,1)$, which EX also checks exactly: $T(\sigma y)=DT(y)D'$ with
$D=\operatorname{diag}(1,s_1,s_2)$, $D'=\operatorname{diag}(1,s_1s_2,1)$ permutes the twelve $t_{ji}$ inside each cut and the sixteen
planes, so $G_\Pi\circ\sigma=G_{\pi(\Pi)}$ and only the $z_4$ component of $\sum_\Pi g_\Pi$ can survive; the $z_4$ component vanishes in EX
(it is the derivative of the wall function $G_w$ at $\pi/4$, zero by the reflection symmetry of Lemma~\ref{lem:fqe-p1-wall}).  By the sign symmetry the tied
gradients have the sign form $(s_1a,s_2a,-s_1s_2c,0,-s_2d,-s_1d)$ of the remark below, so their rank is exactly three
(the characters $s_1,s_2,s_1s_2$ carry independent vectors, as the enclosures give $a,c\ne0$); the bound $-0.05$ for the other planes
and the Hessian enclosures are ball arithmetic (LCC).  A piece $\Omega$ of the sphere is
discharged if, for some $\nu$, $\ell=\sum\nu_\Pi g_\Pi\!\cdot\!\omega\ge0$ and $\ell+\rho q>0$ (or $q>0$) with
$q=\frac12\omega^T[H_\mu]\omega$, evaluated in ball arithmetic. This suffices: on $\Omega$ the average is at least
$t\ell+t^2q=t(\ell+tq)$, and $\ell+tq$ is positive for $0<t\le\rho$, being affine in $t$ with $\ell\ge0$ at $t=0$ and
$\ell+\rho q>0$ at $t=\rho$ (or increasing if $q>0$).  LCC finds such weights on $9012$ cells for the eight points.
\end{proof}

\subsubsection*{The covering}
\begin{proposition}[Interval covering; CCV]\label{prop:fqe-critical-covering}
Let $\mathcal R$ be the part of the chart with $|w|^2\le37/100$, both signs $\mathrm{sg}$.  Every box of the bisection tree is
discharged by one of: (G) a certified positive lower bound of $\delta-C_*d^2$; (HS) $d^2\le1/400$ (then
$\delta\ge\frac3{10}d^2$, strictly greater than $C_*d^2$ when $d>0$, by Theorem~\ref{thm:fqe-orbit-stability}); (core) it lies in a box $\mathcal B_\rho$ of
Lemma~\ref{lem:fqe-critical-local-cone}.
\end{proposition}
\begin{proof}
The starting closed boxes cover $[0,1]^4\times[-5/8,5/8]^2$ in each sign sector, hence the required disc.
A split replaces one box by two closed children whose union is the parent; their common boundary is retained.
Boxes crossing the disc or ordering boundary are retained unless they are wholly outside the domain reduced by symmetry.
Thus complete coverage means that every retained branch ends in one of the three certified alternatives in the statement.
A failed or inconclusive sign test only triggers subdivision; it never certifies a box.

There is no need to choose a unique maximising fidelity plane. Since $G=\max_\Pi G_\Pi$,
any convex combination $G_\mu$ is a lower bound for $G$, including on a tie.
Near equality, the exact contacts and gradient cancellation of the local cone lemma replace a strict sign test at the apex.

On a box with centre $c$ and half-widths $h$ (in $y$) use the weights $\mu$ uniform on the best one to four planes.
(i) $\delta$ is convex in the twelve numbers $t_{ji}$, hence $\delta(y)\ge\delta(c)+\sum_iD_i(t_i(y)-t_i(c))$ with
$D_i=(1+\ln\lambda_i(c))/(12\ln2)$, valid for every point of the box since every chart point is a state, so that $t_i(y)\ge-1$
(the tangent of a convex function holds on its whole domain); a ray whose probability at $c$ is not certified positive is left
out of the tangent and bounded by the maximum of $-x\log_2x$ over its probability box ($-x\log_2x$ increases on $[0,1/e]$).
(ii) $\Phi_\Pi$ is concave in the three entries $m$ of $m_\Pi$, so $d_\Pi^2=2-2\sqrt{\Phi_\Pi}$ is convex there:
$d_\Pi^2(m)\le d_\Pi^2(m_c)+a_\Pi\!\cdot\!(m-m_c)+B$, where the Bregman remainder $B$ is convex in $m-m_c$ and so at most its maximum at the
corners of the box $|m-m_c|\le r$; $r$ is a bound for $|m(y)-m(c)|$.
(iii) $t(y)-t(c)$ and $m(y)-m(c)$ are expanded in $y-c$ to second order, with the exact gradient at $c$ and the interval second
derivatives of the nine entries $T_{\alpha\beta}(y)$ over the box (closed forms, ball arithmetic).
With $g_k=\partial_kG_\mu(c)$ this gives $G\ge G_\mu(c)-\sum_k|g_k|h_k-R$, where $R$ is the sum of the second-order terms of (iii)
inserted into the tangent of (i) and, with factor $C_*$, into the linear term of (ii), and of $C_*\max B$, the first-order terms being exact.
Bisection of the widest side (the floating-point routine \texttt{critical\_float.py} only orders the attempts and certifies nothing)
discharges the chart region $\zeta_1\le\zeta_1'$, $|w|^2\le37/100$ with $208953$ boxes (G) in the sector $\mathrm{sg}=+1$, and
$246044$ (G) and $436$ (HS) for $\mathrm{sg}=-1$, besides $42$ boxes inside the local cone boxes in each sector;
the ordering comparison $\zeta_1>\zeta_1'$ is exact on the dyadic endpoints, while the outside-disc screen has a conservative
positive margin exceeding its floating-point roundoff, so it discards only boxes strictly outside $|w|^2\le37/100$.
The orbit point of the chart ($d=0$, for
$\mathrm{sg}=-1$) is covered by (HS).
\end{proof}

\begin{proof}[Proof of Theorem~\ref{thm:fqe-p1-sharp} for critical states]
If $\tau\ge0.36$, apply Lemma~\ref{lem:fqe-critical-large-tau}.  Otherwise Lemma~\ref{lem:fqe-critical-chart} reduces to
the chart with $|w|^2\le37/100$, which Proposition~\ref{prop:fqe-critical-covering} and Lemma~\ref{lem:fqe-critical-local-cone}
cover: $G>0$, or $d>0$ and $\delta\ge\frac3{10}d^2>C_*d^2$, or $d=0$, except at the equality points $y^{\rm eq}$ where $G=0$ by
\eqref{eq:fqe-p1-cstar}.  The equality points have $w=0$ and are the $144$ points of Theorem~\ref{thm:fqe-p1-sharp}.
\end{proof}

\begin{remark}[Structure at the equality points; numerical picture]\label{rem:fqe-critical-structure}
In the coordinates $z=(\zeta_1,\zeta_2,\zeta_1',1-\zeta_2',w_1,w_2)$ at the equality point with $(\zeta_1,\zeta_2,\zeta_1',\zeta_2')=(0,0,0,1)$, the four tied
plane functions have the gradients $g=(s_1a,\,s_2a,\,-s_1s_2c,\,0,\,-s_2d,\,-s_1d)$, $s_1,s_2=\pm1$, with
$a=0.059039$, $c=0.036251$, $d=0.018765$ (ball enclosures, LCC).  They sum to zero and have rank three, with singular values
$0.1239,0.1239,0.0725$: $G$ has a conical minimum in three directions and is smooth in the other three.  The average of the
four tied Hessians is positive definite on the three-dimensional kernel of the gradients, with eigenvalues
$0.0725$ (the wall direction $\theta$) and $0.7285,0.7285$ (two combinations of the chart coordinates with the $w$-plane).
In particular the dependence on $w$ is first order, in the shape of a square $(\pm1,\pm1)$, so that the loss of
fidelity under moving off $\tau=0$ is compensated at first order by the cone and costs nothing along the smooth directions.
Numerically (not a proof; numerical evidence, NV-1) the minimum of $\delta-C_*d^2$ over critical states with fixed
$\tau$ is about $0.31\,\tau$ for every $\tau\in(0,1]$ (about $0.30\tau$ at $\tau=0.4$ and $0.5$), attained near the equality points for small $\tau$ and,
at $\tau=1$, at a real state with $d^2=2-\sqrt2$.
\end{remark}

\begin{remark}[Non-critical states; numerical evidence only]\label{rem:fqe-critical-noncritical}
The theorem is for critical states; whether $C_*$ is optimal over all states is open.  Evidence (not a proof):
at the equality state, random positive local filters of size $\varepsilon\le0.2$ (the twelve non-critical directions of the
$18$-dimensional LU quotient; $12$ directions per size) give $\delta-C_*d^2\ge3\varepsilon^2$ (NV-2); on states obtained by filtering critical states within $0.4$ of the equality state
(normal-form coordinates) with filters of size $0.05$ to $0.4$ the minimum of $\delta/d^2$ is $0.27$.
These sampled tests found no value below $C_*$. Reduction of a non-critical state to its critical representative does not preserve
$\delta$ (by Section~\ref{sec:fqe-renyi-filter-limit}, one filter on one qubit of a critical state cannot raise any
balanced-cut entropy, but filters on two qubits can raise the mean R\'enyi entropy of order $3/2$), and the general coefficient $k(6+\sqrt6)/12\approx0.05$ is a factor $3.8$ below
$C_*$, so a proof needs a quantitative treatment of the filter directions: near the critical locus a Taylor bound in the filter
variables over this six-parameter chart (a covering of the same size), far from it a stronger distance--entropy comparison.
\end{remark}

